% --- Archon physics formalization source begin ---
% archon:physics
% archon:covers PhyXMiniProblems/problem_phyx_mini_0972.lean
% archon:source-report reports/phyx_mini/problem_phyx_mini_0972.source.json

\chapter{Physics problem phyx\_mini\_0972}
\label{ch:PhyXMiniProblems_problem_phyx_mini_0972}

\paragraph{Problem source.}
\#\# Physical scenario

Blood contains positive and negative ions and thus is a conductor. A blood vessel, therefore, can be viewed as an electrical wire. We can even picture the flowing blood as a series of parallel conducting slabs whose thickness is the diameter \textbackslash{}( d \textbackslash{}) of the vessel moving with speed \textbackslash{}( v \textbackslash{}).

\#\# Question

If you expect that the blood will be flowing at 15\textbackslash{},\textbackslash{}text\{cm/s\} for a vessel 5.0\textbackslash{},\textbackslash{}text\{mm\} in diameter, what strength of magnetic field will you need to produce a potential difference of 1.0\textbackslash{},\textbackslash{}text\{mV\}?

\#\# Answer choices

- A: A: 13.2T

- B: B: 1.32T

- C: C: 0.013T

- D: D: 0.00013T

\#\# Figure caption (auxiliary; use the image as primary evidence)

The image shows a cylindrical rod moving through a magnetic field. Here are the details:

- The cylinder is depicted in red and is oriented horizontally.
- Blue "X" symbols indicate the direction of the magnetic field \textbackslash{}(\textbackslash{}vec\{B\}\textbackslash{}), which is directed into the page.
- An arrow labeled \textbackslash{}(v\textbackslash{}) in green indicates the velocity of the cylinder, pointing to the right.
- The diameter of the cylinder is labeled as \textbackslash{}(d\textbackslash{}).
- The magnetic field \textbackslash{}(\textbackslash{}vec\{B\}\textbackslash{}) is represented by a vector above the cylinder, with the symbol \textbackslash{}(\textbackslash{}vec\{B\}\textbackslash{}) placed near it.

The scene illustrates the interaction of the cylinder with the magnetic field as it moves horizontally.

\#\# Dataset metadata

Electromagnetic Induction; Physical Model Grounding Reasoning, Multi-Formula Reasoning

\paragraph{Recorded answer/context.}
B: B: 1.32T

\paragraph{Figure/image path.}
/root/proposal\_for\_physic/science-mango/phyx\_mini\_run/phyx\_data/test\_image/972.png

\paragraph{Formalization target.}
create a compiling Lean file with sorry bodies at `PhyXMiniProblems/problem\_phyx\_mini\_0972.lean`. The Lean declarations must preserve the physical quantities, dimensions or dimensional roles, figure labels, governing-law hypotheses, and final relation expressed by this problem.
Use Mathlib/Physlib names found through LeanExplore where available. If a physics API is missing, introduce faithful local abstractions rather than scalar placeholder aliases.

\begin{theorem}[Physics formalization target]
\label{thm:physics:phyx_mini_0972:target}
\lean{PhyXMiniProblems.ProblemPhyXMini0972.problem_phyx_mini_0972}
\uses{def:physics:phyx-mini-0972:phyxminiproblems-problemphyxmini0972-magneticfluxdensityinteslas, def:physics:phyx-mini-0972:phyxminiproblems-problemphyxmini0972-bloodvesselmotionalemfsetup, def:physics:phyx-mini-0972:phyxminiproblems-problemphyxmini0972-matcheswrittenbloodvesselscenario, def:physics:phyx-mini-0972:phyxminiproblems-problemphyxmini0972-matchessuppliedbloodvesselfigure, def:physics:phyx-mini-0972:phyxminiproblems-problemphyxmini0972-hasgivenbloodvesselmeasurements, def:physics:phyx-mini-0972:phyxminiproblems-problemphyxmini0972-satisfiesperpendicularmotionalemflaw, def:physics:phyx-mini-0972:phyxminiproblems-problemphyxmini0972-recordeddatasetanswer, def:physics:phyx-mini-0972:phyxminiproblems-problemphyxmini0972-isclosestdisplayedfieldstrength}
The assigned autoformalize agent should translate this physics problem into Lean declarations in the covered file, with theorem and lemma proof bodies written as `by sorry`.
\end{theorem}
\begin{proof}
This is an autoformalization task, not a proof task. Produce statements that can later be proved by the physics prover without weakening the physical model.
\end{proof}

% --- Archon declaration topology begin ---
\section{Declaration topology}
\begin{definition}[magnetic Flux Density Dimension]
  \label{def:physics:phyx-mini-0972:phyxminiproblems-problemphyxmini0972-magneticfluxdensitydimension}
  \lean{PhyXMiniProblems.ProblemPhyXMini0972.magneticFluxDensityDimension}
  Magnetic-flux density, measured in teslas in coherent SI, has dimension M T⁻¹ C⁻¹
\end{definition}
\begin{proof}
  This is the declared model definition of magnetic Flux Density Dimension.
\end{proof}
\begin{definition}[potential Difference Dimension]
  \label{def:physics:phyx-mini-0972:phyxminiproblems-problemphyxmini0972-potentialdifferencedimension}
  \lean{PhyXMiniProblems.ProblemPhyXMini0972.potentialDifferenceDimension}
  \uses{def:physics:phyx-mini-0972:phyxminiproblems-problemphyxmini0972-magneticfluxdensitydimension}
  Electric potential difference has dimension B · velocity · length
\end{definition}
\begin{proof}
  This is the declared model definition of potential Difference Dimension.
\end{proof}
\begin{definition}[Length Magnitude Quantity]
  \label{def:physics:phyx-mini-0972:phyxminiproblems-problemphyxmini0972-lengthmagnitudequantity}
  \lean{PhyXMiniProblems.ProblemPhyXMini0972.LengthMagnitudeQuantity}
  A nonnegative, unit-independent physical length
\end{definition}
\begin{proof}
  This is the declared model definition of Length Magnitude Quantity.
\end{proof}
\begin{definition}[Speed Magnitude Quantity]
  \label{def:physics:phyx-mini-0972:phyxminiproblems-problemphyxmini0972-speedmagnitudequantity}
  \lean{PhyXMiniProblems.ProblemPhyXMini0972.SpeedMagnitudeQuantity}
  A nonnegative, unit-independent physical speed
\end{definition}
\begin{proof}
  This is the declared model definition of Speed Magnitude Quantity.
\end{proof}
\begin{definition}[Magnetic Flux Density Magnitude Quantity]
  \label{def:physics:phyx-mini-0972:phyxminiproblems-problemphyxmini0972-magneticfluxdensitymagnitudequantity}
  \lean{PhyXMiniProblems.ProblemPhyXMini0972.MagneticFluxDensityMagnitudeQuantity}
  \uses{def:physics:phyx-mini-0972:phyxminiproblems-problemphyxmini0972-magneticfluxdensitydimension}
  A nonnegative magnetic-flux-density magnitude
\end{definition}
\begin{proof}
  This is the declared model definition of Magnetic Flux Density Magnitude Quantity.
\end{proof}
\begin{definition}[Potential Difference Magnitude Quantity]
  \label{def:physics:phyx-mini-0972:phyxminiproblems-problemphyxmini0972-potentialdifferencemagnitudequantity}
  \lean{PhyXMiniProblems.ProblemPhyXMini0972.PotentialDifferenceMagnitudeQuantity}
  \uses{def:physics:phyx-mini-0972:phyxminiproblems-problemphyxmini0972-potentialdifferencedimension}
  A nonnegative electric-potential-difference magnitude
\end{definition}
\begin{proof}
  This is the declared model definition of Potential Difference Magnitude Quantity.
\end{proof}
\begin{definition}[coherent SIReadout]
  \label{def:physics:phyx-mini-0972:phyxminiproblems-problemphyxmini0972-coherentsireadout}
  \lean{PhyXMiniProblems.ProblemPhyXMini0972.coherentSIReadout}
  Coherent-SI readout of a nonnegative dimensionful quantity
\end{definition}
\begin{proof}
  This is the declared model definition of coherent SIReadout.
\end{proof}
\begin{definition}[length In Meters]
  \label{def:physics:phyx-mini-0972:phyxminiproblems-problemphyxmini0972-lengthinmeters}
  \lean{PhyXMiniProblems.ProblemPhyXMini0972.lengthInMeters}
  \uses{def:physics:phyx-mini-0972:phyxminiproblems-problemphyxmini0972-lengthmagnitudequantity, def:physics:phyx-mini-0972:phyxminiproblems-problemphyxmini0972-coherentsireadout}
  Read a length in metres
\end{definition}
\begin{proof}
  This is the declared model definition of length In Meters.
\end{proof}
\begin{definition}[length In Millimeters]
  \label{def:physics:phyx-mini-0972:phyxminiproblems-problemphyxmini0972-lengthinmillimeters}
  \lean{PhyXMiniProblems.ProblemPhyXMini0972.lengthInMillimeters}
  \uses{def:physics:phyx-mini-0972:phyxminiproblems-problemphyxmini0972-lengthmagnitudequantity}
  Read a length in millimetres
\end{definition}
\begin{proof}
  This is the declared model definition of length In Millimeters.
\end{proof}
\begin{definition}[speed In Meters Per Second]
  \label{def:physics:phyx-mini-0972:phyxminiproblems-problemphyxmini0972-speedinmeterspersecond}
  \lean{PhyXMiniProblems.ProblemPhyXMini0972.speedInMetersPerSecond}
  \uses{def:physics:phyx-mini-0972:phyxminiproblems-problemphyxmini0972-speedmagnitudequantity, def:physics:phyx-mini-0972:phyxminiproblems-problemphyxmini0972-coherentsireadout}
  Read a speed in metres per second
\end{definition}
\begin{proof}
  This is the declared model definition of speed In Meters Per Second.
\end{proof}
\begin{definition}[speed In Centimeters Per Second]
  \label{def:physics:phyx-mini-0972:phyxminiproblems-problemphyxmini0972-speedincentimeterspersecond}
  \lean{PhyXMiniProblems.ProblemPhyXMini0972.speedInCentimetersPerSecond}
  \uses{def:physics:phyx-mini-0972:phyxminiproblems-problemphyxmini0972-speedmagnitudequantity}
  Read a speed in centimetres per second
\end{definition}
\begin{proof}
  This is the declared model definition of speed In Centimeters Per Second.
\end{proof}
\begin{definition}[magnetic Flux Density In Teslas]
  \label{def:physics:phyx-mini-0972:phyxminiproblems-problemphyxmini0972-magneticfluxdensityinteslas}
  \lean{PhyXMiniProblems.ProblemPhyXMini0972.magneticFluxDensityInTeslas}
  \uses{def:physics:phyx-mini-0972:phyxminiproblems-problemphyxmini0972-magneticfluxdensitymagnitudequantity, def:physics:phyx-mini-0972:phyxminiproblems-problemphyxmini0972-coherentsireadout}
  Read a magnetic-flux-density magnitude in teslas
\end{definition}
\begin{proof}
  This is the declared model definition of magnetic Flux Density In Teslas.
\end{proof}
\begin{definition}[potential Difference In Volts]
  \label{def:physics:phyx-mini-0972:phyxminiproblems-problemphyxmini0972-potentialdifferenceinvolts}
  \lean{PhyXMiniProblems.ProblemPhyXMini0972.potentialDifferenceInVolts}
  \uses{def:physics:phyx-mini-0972:phyxminiproblems-problemphyxmini0972-potentialdifferencemagnitudequantity, def:physics:phyx-mini-0972:phyxminiproblems-problemphyxmini0972-coherentsireadout}
  Read a potential-difference magnitude in volts
\end{definition}
\begin{proof}
  This is the declared model definition of potential Difference In Volts.
\end{proof}
\begin{definition}[potential Difference In Millivolts]
  \label{def:physics:phyx-mini-0972:phyxminiproblems-problemphyxmini0972-potentialdifferenceinmillivolts}
  \lean{PhyXMiniProblems.ProblemPhyXMini0972.potentialDifferenceInMillivolts}
  \uses{def:physics:phyx-mini-0972:phyxminiproblems-problemphyxmini0972-potentialdifferencemagnitudequantity, def:physics:phyx-mini-0972:phyxminiproblems-problemphyxmini0972-potentialdifferenceinvolts}
  Read a potential-difference magnitude in millivolts
\end{definition}
\begin{proof}
  This is the declared model definition of potential Difference In Millivolts.
\end{proof}
\begin{definition}[Charge Sign]
  \label{def:physics:phyx-mini-0972:phyxminiproblems-problemphyxmini0972-chargesign}
  \lean{PhyXMiniProblems.ProblemPhyXMini0972.ChargeSign}
  The two signs of mobile ionic charge carriers mentioned in the scenario
\end{definition}
\begin{proof}
  This is the declared model definition of Charge Sign.
\end{proof}
\begin{definition}[Conducting Blood]
  \label{def:physics:phyx-mini-0972:phyxminiproblems-problemphyxmini0972-conductingblood}
  \lean{PhyXMiniProblems.ProblemPhyXMini0972.ConductingBlood}
  \uses{def:physics:phyx-mini-0972:phyxminiproblems-problemphyxmini0972-chargesign}
  An abstract conducting-blood model retaining mobile positive and negative ion species without reducing either charge or conductivity to a scalar alias
\end{definition}
\begin{proof}
  This is the declared model definition of Conducting Blood.
\end{proof}
\begin{definition}[Vessel Electrical Model]
  \label{def:physics:phyx-mini-0972:phyxminiproblems-problemphyxmini0972-vesselelectricalmodel}
  \lean{PhyXMiniProblems.ProblemPhyXMini0972.VesselElectricalModel}
  Electrical idealization assigned to the vessel and its contents
\end{definition}
\begin{proof}
  This is the declared model definition of Vessel Electrical Model.
\end{proof}
\begin{definition}[Blood Flow Idealization]
  \label{def:physics:phyx-mini-0972:phyxminiproblems-problemphyxmini0972-bloodflowidealization}
  \lean{PhyXMiniProblems.ProblemPhyXMini0972.BloodFlowIdealization}
  Kinematic idealization used for the flowing blood
\end{definition}
\begin{proof}
  This is the declared model definition of Blood Flow Idealization.
\end{proof}
\begin{definition}[Planar Direction]
  \label{def:physics:phyx-mini-0972:phyxminiproblems-problemphyxmini0972-planardirection}
  \lean{PhyXMiniProblems.ProblemPhyXMini0972.PlanarDirection}
  Directions distinguished in the plane of the supplied image
\end{definition}
\begin{proof}
  This is the declared model definition of Planar Direction.
\end{proof}
\begin{definition}[Planar Axis]
  \label{def:physics:phyx-mini-0972:phyxminiproblems-problemphyxmini0972-planaraxis}
  \lean{PhyXMiniProblems.ProblemPhyXMini0972.PlanarAxis}
  Axes distinguished in the plane of the supplied image
\end{definition}
\begin{proof}
  This is the declared model definition of Planar Axis.
\end{proof}
\begin{definition}[Page Normal Direction]
  \label{def:physics:phyx-mini-0972:phyxminiproblems-problemphyxmini0972-pagenormaldirection}
  \lean{PhyXMiniProblems.ProblemPhyXMini0972.PageNormalDirection}
  Directions normal to the page of the supplied image
\end{definition}
\begin{proof}
  This is the declared model definition of Page Normal Direction.
\end{proof}
\begin{definition}[Magnetic Field Glyph]
  \label{def:physics:phyx-mini-0972:phyxminiproblems-problemphyxmini0972-magneticfieldglyph}
  \lean{PhyXMiniProblems.ProblemPhyXMini0972.MagneticFieldGlyph}
  Glyph convention for a magnetic field normal to the page
\end{definition}
\begin{proof}
  This is the declared model definition of Magnetic Field Glyph.
\end{proof}
\begin{definition}[Relative Orientation]
  \label{def:physics:phyx-mini-0972:phyxminiproblems-problemphyxmini0972-relativeorientation}
  \lean{PhyXMiniProblems.ProblemPhyXMini0972.RelativeOrientation}
  Relative orientation of two physical directions
\end{definition}
\begin{proof}
  This is the declared model definition of Relative Orientation.
\end{proof}
\begin{definition}[Blood Vessel Figure]
  \label{def:physics:phyx-mini-0972:phyxminiproblems-problemphyxmini0972-bloodvesselfigure}
  \lean{PhyXMiniProblems.ProblemPhyXMini0972.BloodVesselFigure}
  \uses{def:physics:phyx-mini-0972:phyxminiproblems-problemphyxmini0972-planardirection, def:physics:phyx-mini-0972:phyxminiproblems-problemphyxmini0972-planaraxis, def:physics:phyx-mini-0972:phyxminiproblems-problemphyxmini0972-magneticfieldglyph}
  Literal presentation data transcribed from the primary raster
\end{definition}
\begin{proof}
  This is the declared model definition of Blood Vessel Figure.
\end{proof}
\begin{definition}[Blood Vessel Motional Emf Setup]
  \label{def:physics:phyx-mini-0972:phyxminiproblems-problemphyxmini0972-bloodvesselmotionalemfsetup}
  \lean{PhyXMiniProblems.ProblemPhyXMini0972.BloodVesselMotionalEmfSetup}
  \uses{def:physics:phyx-mini-0972:phyxminiproblems-problemphyxmini0972-lengthmagnitudequantity, def:physics:phyx-mini-0972:phyxminiproblems-problemphyxmini0972-speedmagnitudequantity, def:physics:phyx-mini-0972:phyxminiproblems-problemphyxmini0972-magneticfluxdensitymagnitudequantity, def:physics:phyx-mini-0972:phyxminiproblems-problemphyxmini0972-potentialdifferencemagnitudequantity, def:physics:phyx-mini-0972:phyxminiproblems-problemphyxmini0972-conductingblood, def:physics:phyx-mini-0972:phyxminiproblems-problemphyxmini0972-vesselelectricalmodel, def:physics:phyx-mini-0972:phyxminiproblems-problemphyxmini0972-bloodflowidealization, def:physics:phyx-mini-0972:phyxminiproblems-problemphyxmini0972-planardirection, def:physics:phyx-mini-0972:phyxminiproblems-problemphyxmini0972-planaraxis, def:physics:phyx-mini-0972:phyxminiproblems-problemphyxmini0972-pagenormaldirection, def:physics:phyx-mini-0972:phyxminiproblems-problemphyxmini0972-relativeorientation, def:physics:phyx-mini-0972:phyxminiproblems-problemphyxmini0972-bloodvesselfigure}
  Independent physical observables for the vessel experiment. The slab thickness and vessel diameter are kept distinct until the written model equates them. Likewise, the requested field strength is not computed by this record
\end{definition}
\begin{proof}
  This is the declared model definition of Blood Vessel Motional Emf Setup.
\end{proof}
\begin{definition}[Matches Written Blood Vessel Scenario]
  \label{def:physics:phyx-mini-0972:phyxminiproblems-problemphyxmini0972-matcheswrittenbloodvesselscenario}
  \lean{PhyXMiniProblems.ProblemPhyXMini0972.MatchesWrittenBloodVesselScenario}
  \uses{def:physics:phyx-mini-0972:phyxminiproblems-problemphyxmini0972-bloodvesselmotionalemfsetup}
  The conductor-and-moving-slab idealization stated in the written scenario
\end{definition}
\begin{proof}
  This is the declared model definition of Matches Written Blood Vessel Scenario.
\end{proof}
\begin{definition}[Matches Supplied Blood Vessel Figure]
  \label{def:physics:phyx-mini-0972:phyxminiproblems-problemphyxmini0972-matchessuppliedbloodvesselfigure}
  \lean{PhyXMiniProblems.ProblemPhyXMini0972.MatchesSuppliedBloodVesselFigure}
  \uses{def:physics:phyx-mini-0972:phyxminiproblems-problemphyxmini0972-bloodvesselmotionalemfsetup}
  Literal figure evidence and its calibration to the physical setup. The standard cross-glyph convention is recorded explicitly as “into the page.”
\end{definition}
\begin{proof}
  This is the declared model definition of Matches Supplied Blood Vessel Figure.
\end{proof}
\begin{definition}[Has Given Blood Vessel Measurements]
  \label{def:physics:phyx-mini-0972:phyxminiproblems-problemphyxmini0972-hasgivenbloodvesselmeasurements}
  \lean{PhyXMiniProblems.ProblemPhyXMini0972.HasGivenBloodVesselMeasurements}
  \uses{def:physics:phyx-mini-0972:phyxminiproblems-problemphyxmini0972-lengthinmillimeters, def:physics:phyx-mini-0972:phyxminiproblems-problemphyxmini0972-speedincentimeterspersecond, def:physics:phyx-mini-0972:phyxminiproblems-problemphyxmini0972-potentialdifferenceinmillivolts, def:physics:phyx-mini-0972:phyxminiproblems-problemphyxmini0972-bloodvesselmotionalemfsetup}
  The three nominal measurements printed in the question. No magnetic-field value occurs here: it remains the quantity requested by the problem
\end{definition}
\begin{proof}
  This is the declared model definition of Has Given Blood Vessel Measurements.
\end{proof}
\begin{definition}[Satisfies Perpendicular Motional Emf Law]
  \label{def:physics:phyx-mini-0972:phyxminiproblems-problemphyxmini0972-satisfiesperpendicularmotionalemflaw}
  \lean{PhyXMiniProblems.ProblemPhyXMini0972.SatisfiesPerpendicularMotionalEmfLaw}
  \uses{def:physics:phyx-mini-0972:phyxminiproblems-problemphyxmini0972-lengthinmeters, def:physics:phyx-mini-0972:phyxminiproblems-problemphyxmini0972-speedinmeterspersecond, def:physics:phyx-mini-0972:phyxminiproblems-problemphyxmini0972-magneticfluxdensityinteslas, def:physics:phyx-mini-0972:phyxminiproblems-problemphyxmini0972-potentialdifferenceinvolts, def:physics:phyx-mini-0972:phyxminiproblems-problemphyxmini0972-bloodvesselmotionalemfsetup}
  For mutually perpendicular B, v, and charge-separation direction, the motional-emf magnitude law is ΔV = B v d. This premise is a governing law, not the requested numerical field strength
\end{definition}
\begin{proof}
  This is the declared model definition of Satisfies Perpendicular Motional Emf Law.
\end{proof}
\begin{lemma}[magnetic Field Strength eq potential Difference div speed mul thickness]
  \label{lem:physics:phyx-mini-0972:phyxminiproblems-problemphyxmini0972-magneticfieldstrength-eq-potentialdifference-div-speed-mul-thickness}
  \lean{PhyXMiniProblems.ProblemPhyXMini0972.magneticFieldStrength_eq_potentialDifference_div_speed_mul_thickness}
  \uses{def:physics:phyx-mini-0972:phyxminiproblems-problemphyxmini0972-lengthinmeters, def:physics:phyx-mini-0972:phyxminiproblems-problemphyxmini0972-speedinmeterspersecond, def:physics:phyx-mini-0972:phyxminiproblems-problemphyxmini0972-magneticfluxdensityinteslas, def:physics:phyx-mini-0972:phyxminiproblems-problemphyxmini0972-potentialdifferenceinvolts, def:physics:phyx-mini-0972:phyxminiproblems-problemphyxmini0972-bloodvesselmotionalemfsetup, def:physics:phyx-mini-0972:phyxminiproblems-problemphyxmini0972-satisfiesperpendicularmotionalemflaw}
  Solving the motional-emf law for a nondegenerate magnetic field strength
\end{lemma}
\begin{proof}
  The conclusion follows from the governing relations and figure readouts named in the statement of magnetic Field Strength eq potential Difference div speed mul thickness.
\end{proof}
\begin{definition}[Answer Choice]
  \label{def:physics:phyx-mini-0972:phyxminiproblems-problemphyxmini0972-answerchoice}
  \lean{PhyXMiniProblems.ProblemPhyXMini0972.AnswerChoice}
  Labels of the four field strengths displayed in the source
\end{definition}
\begin{proof}
  This is the declared model definition of Answer Choice.
\end{proof}
\begin{definition}[displayed Field Strength In Teslas]
  \label{def:physics:phyx-mini-0972:phyxminiproblems-problemphyxmini0972-answerchoice-displayedfieldstrengthinteslas}
  \lean{PhyXMiniProblems.ProblemPhyXMini0972.AnswerChoice.displayedFieldStrengthInTeslas}
  \uses{def:physics:phyx-mini-0972:phyxminiproblems-problemphyxmini0972-answerchoice}
  Magnetic-field strength in teslas displayed beside an answer label
\end{definition}
\begin{proof}
  This is the declared model definition of displayed Field Strength In Teslas.
\end{proof}
\begin{definition}[recorded Dataset Answer]
  \label{def:physics:phyx-mini-0972:phyxminiproblems-problemphyxmini0972-recordeddatasetanswer}
  \lean{PhyXMiniProblems.ProblemPhyXMini0972.recordedDatasetAnswer}
  \uses{def:physics:phyx-mini-0972:phyxminiproblems-problemphyxmini0972-answerchoice}
  The dataset records answer label B
\end{definition}
\begin{proof}
  This is the declared model definition of recorded Dataset Answer.
\end{proof}
\begin{definition}[Is Closest Displayed Field Strength]
  \label{def:physics:phyx-mini-0972:phyxminiproblems-problemphyxmini0972-isclosestdisplayedfieldstrength}
  \lean{PhyXMiniProblems.ProblemPhyXMini0972.IsClosestDisplayedFieldStrength}
  \uses{def:physics:phyx-mini-0972:phyxminiproblems-problemphyxmini0972-magneticfluxdensityinteslas, def:physics:phyx-mini-0972:phyxminiproblems-problemphyxmini0972-bloodvesselmotionalemfsetup, def:physics:phyx-mini-0972:phyxminiproblems-problemphyxmini0972-answerchoice}
  A choice is strictly closer than every alternative to the required field
\end{definition}
\begin{proof}
  This is the declared model definition of Is Closest Displayed Field Strength.
\end{proof}
% --- Archon declaration topology end ---
% --- Archon physics formalization source end ---
